% ═══════════════════════════════════════════════════════════════
% LERNBLATT: HAY (Repaso) + ESTAR
% Spanisch - Greenhouse Schools Rostock
% ═══════════════════════════════════════════════════════════════
% Thema: HAY vs ESTAR
% Klasse 7 | Unidad 2 | Mi instituto | DS 02
% ═══════════════════════════════════════════════════════════════

\documentclass[11pt, a4paper]{article}

\usepackage[utf8]{inputenc}
\usepackage[T1]{fontenc}
\usepackage[ngerman]{babel}

% --- SCHRIFTEN ---
\usepackage{helvet}
\renewcommand{\familydefault}{\sfdefault}
\usepackage{palatino}
\newcommand{\seriftext}[1]{{\fontfamily{ppl}\selectfont #1}}

\usepackage{microtype}
\usepackage[margin=1.8cm, top=2cm, bottom=1.5cm]{geometry}
\usepackage{enumitem}
\usepackage{tabularx}
\usepackage{booktabs}
\usepackage{array}
\usepackage{multicol}
\usepackage{tikz}

\usepackage{fancyhdr}
\usepackage{lastpage}

% ═══════════════════════════════════════════════════════════════
% VARIABLEN
% ═══════════════════════════════════════════════════════════════

\newcommand{\lektion}{2}
\newcommand{\thema}{Hay y estar}
\newcommand{\untertitel}{Wo ist was? Was gibt es?}
\newcommand{\klasse}{7}
\newcommand{\lernziele}{hay (Wiederholung) · estar (Konjugation) · Präpositionen · hay vs estar}

% ═══════════════════════════════════════════════════════════════
% FARBEN
% ═══════════════════════════════════════════════════════════════

\usepackage{xcolor}
\usepackage{amssymb}

\definecolor{colorrojo}{HTML}{C41E3A}       % Spanisch-Akzent
\definecolor{colorazul}{HTML}{1565C0}
\definecolor{colorverde}{HTML}{2E7D32}
\definecolor{colorgris}{HTML}{757575}
\definecolor{colornegro}{HTML}{222222}

\definecolor{hellgrau}{HTML}{F5F5F5}
\definecolor{rahmen}{HTML}{CCCCCC}
\definecolor{akzent}{HTML}{666666}

% ═══════════════════════════════════════════════════════════════
% BOXEN
% ═══════════════════════════════════════════════════════════════

\usepackage{tcolorbox}
\tcbuselibrary{skins}

% Lernziel-Box
\newtcolorbox{lernzielbox}{
    colback=hellgrau, colframe=hellgrau, boxrule=0pt, arc=0pt,
    left=10pt, right=10pt, top=8pt, bottom=8pt,
    before skip=6pt, after skip=8pt
}

% Grammatik-Box (Randlinie links)
\newtcolorbox{grammatikbox}[1][]{
    colback=white, colframe=white, boxrule=0pt, arc=0pt,
    left=10pt, right=6pt, top=6pt, bottom=6pt,
    borderline west={3pt}{0pt}{akzent},
    before skip=4pt, after skip=4pt
}

% Merk-Box (Linie oben)
\newtcolorbox{merkbox}[1][]{
    colback=hellgrau, colframe=hellgrau, boxrule=0pt, arc=0pt,
    left=10pt, right=10pt, top=8pt, bottom=8pt,
    borderline north={2pt}{0pt}{colornegro},
    before skip=6pt, after skip=6pt
}

% Vokabel-Box
\newtcolorbox{vokabelbox}[1][]{
    colback=white, colframe=rahmen, boxrule=0.5pt, arc=0pt,
    left=8pt, right=8pt, top=6pt, bottom=6pt,
    before skip=4pt, after skip=4pt
}

% Tipp-Box
\newtcolorbox{tippbox}{
    colback=white, colframe=white, boxrule=0pt, arc=0pt,
    left=8pt, right=4pt, top=3pt, bottom=3pt,
    borderline west={2pt}{0pt}{akzent}
}

% Fehler-Box
\newtcolorbox{fehlerbox}{
    colback=white, colframe=colorrojo, boxrule=0.5pt, arc=0pt,
    left=8pt, right=8pt, top=6pt, bottom=6pt
}

% Vergleichs-Box
\newtcolorbox{vergleichbox}{
    colback=hellgrau, colframe=hellgrau, boxrule=0pt, arc=0pt,
    left=10pt, right=10pt, top=8pt, bottom=8pt,
    borderline north={2pt}{0pt}{colorrojo},
    before skip=6pt, after skip=6pt
}

% ═══════════════════════════════════════════════════════════════
% BEFEHLE
% ═══════════════════════════════════════════════════════════════

% Spanisch: Serif + fett
\newcommand{\es}[1]{\seriftext{\textbf{#1}}}

% Deutsch: kursiv + grau
\newcommand{\de}[1]{\textit{\textcolor{akzent}{#1}}}

% Grammatikbegriff: Kapitälchen
\newcommand{\gram}[1]{\textsc{#1}}

% AB-Verweis
\newcommand{\abmark}[1]{{\footnotesize\textcolor{akzent}{$\rightarrow$ AB #1}}}

% Abschnittstitel
\newcommand{\abschnitt}[2]{\noindent\textbf{\large #1. #2}}

% Stammwechsel hervorheben
\newcommand{\stamm}[1]{\textcolor{colorrojo}{#1}}

\setlength{\parindent}{0pt}
\setlength{\parskip}{5pt}

% ═══════════════════════════════════════════════════════════════
% KOPF- UND FUSSZEILE
% ═══════════════════════════════════════════════════════════════

\pagestyle{fancy}
\fancyhf{}
\renewcommand{\headrulewidth}{0.5pt}
\renewcommand{\footrulewidth}{0pt}
\lhead{\footnotesize\textsc{Spanisch} · Klasse \klasse}
\rhead{\footnotesize\thema}
\cfoot{\footnotesize\thepage\,/\,\pageref{LastPage}}

% ═══════════════════════════════════════════════════════════════
% DOKUMENT
% ═══════════════════════════════════════════════════════════════

\begin{document}

% --- TITEL ---
\begin{center}
    {\LARGE\bfseries \thema}\\[3pt]
    {\small \untertitel}
\end{center}

\vspace{4pt}

% --- EINLEITUNG ---
Du kennst schon \es{hay} -- damit sagst du, \textbf{was es gibt}. Aber wie sagst du, \textbf{wo etwas ist}? Dafür brauchst du das Verb \es{estar}!

\vspace{2pt}

% --- LERNZIELE ---
\begin{lernzielbox}
\textbf{Das lernst du:}\quad \lernziele
\end{lernzielbox}

% ═══════════════════════════════════════════════════════════════
% ZWEI SPALTEN
% ═══════════════════════════════════════════════════════════════

\begin{multicols}{2}

% ---------------------------------------------------------------
% ABSCHNITT 1: HAY - WIEDERHOLUNG
% ---------------------------------------------------------------
\abschnitt{1}{Hay} \de{-- Es gibt (Repaso)}

\vspace{4pt}

\begin{grammatikbox}
\textbf{\gram{hay}} = es gibt

\vspace{4pt}

\textbf{Unveränderlich} -- Singular \& Plural gleich!

\vspace{6pt}

{\small
\begin{tabular}{@{}p{3.8cm}l@{}}
\textbf{+ artículo indefinido} & \es{Hay \textbf{un} libro.} \\[0.3em]
\textbf{+ número} & \es{Hay \textbf{tres} sillas.} \\[0.3em]
\textbf{+ palabra de cantidad} & \es{Hay \textbf{muchos} libros.} \\
 & \es{Hay \textbf{pocos} cuadernos.} \\[0.3em]
\textbf{+ sin artículo} & \es{Hay \textbf{libros}.} \\
 & \es{¿Hay \textbf{leche}?} \\
\end{tabular}}
\end{grammatikbox}

\vspace{4pt}

\begin{tippbox}
\footnotesize\textit{Neu:} Ohne Artikel = allgemein, Mehrzahl oder nicht zählbar.\\
\es{Hay libros.} \de{Es gibt Bücher (irgendwelche).}
\end{tippbox}

\vspace{4pt}

\begin{fehlerbox}
{\small\textbf{FALSCH:} \textcolor{colorrojo}{*Hay \textit{el} libro.} / \textcolor{colorrojo}{*Hay \textit{la} mesa.}

\vspace{2pt}

\textbf{Merke:} Nach \es{hay} kommt \textbf{nie} el/la!}
\end{fehlerbox}

\abmark{1, 2}

\columnbreak

% ---------------------------------------------------------------
% ABSCHNITT 2: ESTAR - NEU
% ---------------------------------------------------------------
\abschnitt{2}{Estar} \de{-- Sein (Ort/Zustand)}

\vspace{4pt}

\begin{grammatikbox}
\textbf{\gram{estar}} = sein (für \textbf{Ort/Zustand})

\vspace{6pt}

{\small
\begin{tabular}{@{}ll@{}}
yo & \es{estoy} \\
tú & \es{estás} \\
él/ella/usted & \es{está} \\
nosotros/as & \es{estamos} \\
vosotros/as & \es{estáis} \\
ellos/ellas/ustedes & \es{están} \\
\end{tabular}}

\vspace{6pt}

{\footnotesize\textit{Merke:} Alle Formen außer \es{estamos} haben Akzent!}
\end{grammatikbox}

\vspace{4pt}

\textbf{Beispiele:}\\
{\small
\es{¿Dónde está el libro?} \de{Wo ist das Buch?}\\
\es{El libro está en la mesa.} \de{Das Buch ist auf dem Tisch.}}

\abmark{3}

\vspace{8pt}

% ---------------------------------------------------------------
% ABSCHNITT 3: PRÄPOSITIONEN
% ---------------------------------------------------------------
\abschnitt{3}{Preposiciones} \de{-- Präpositionen}

\vspace{4pt}

\begin{vokabelbox}
{\small
\begin{tabular}{@{}ll@{}}
\es{en} & \de{in, auf, an} \\[0.3em]
\es{encima de} & \de{auf, oberhalb von} \\[0.3em]
\es{debajo de} & \de{unter} \\[0.3em]
\es{al lado de} & \de{neben} \\[0.3em]
\es{delante de} & \de{vor} \\[0.3em]
\es{detrás de} & \de{hinter} \\[0.3em]
\es{entre} & \de{zwischen} \\
\end{tabular}}
\end{vokabelbox}

\abmark{3}

\end{multicols}

% ═══════════════════════════════════════════════════════════════
% ABSCHNITT 4: HAY VS ESTAR (volle Breite)
% ═══════════════════════════════════════════════════════════════

\vspace{6pt}

\abschnitt{4}{Hay vs Estar} \de{-- Der Unterschied}

\vspace{4pt}

\begin{vergleichbox}
\begin{multicols}{2}

{\large\bfseries hay} = \textbf{Existenz}

\vspace{4pt}

\textit{Gibt es das?} (unbekannt/unbestimmt)

\vspace{4pt}

{\small
\es{Hay \textbf{un} libro en la mesa.}\\
\de{Es gibt \textbf{ein} Buch auf dem Tisch.}\\
\de{(irgendein Buch -- ich weiß nicht welches)}}

\vspace{4pt}

$\rightarrow$ \textbf{un/una}, Zahlen, muchos, sin artículo

\columnbreak

{\large\bfseries estar} = \textbf{Ort}

\vspace{4pt}

\textit{Wo ist es?} (bekannt/bestimmt)

\vspace{4pt}

{\small
\es{\textbf{El} libro está en la mesa.}\\
\de{\textbf{Das} Buch ist auf dem Tisch.}\\
\de{(ein bestimmtes Buch -- ich weiß welches)}}

\vspace{4pt}

$\rightarrow$ \textbf{el/la}, mi/tu/su, este/ese

\end{multicols}
\end{vergleichbox}

\vspace{6pt}

\begin{multicols}{2}

\textbf{Weitere Beispiele:}

{\small
\begin{tabular}{@{}p{5.8cm}l@{}}
\es{¿Hay una goma?} & \de{Gibt es einen Radierer?} \\[0.3em]
\es{Sí, hay una goma.} & \de{Ja, es gibt einen.} \\[0.3em]
\es{¿Dónde está la goma?} & \de{Wo ist der Radierer?} \\[0.3em]
\es{La goma está en el estuche.} & \de{Der Radierer ist in der Federtasche.} \\
\end{tabular}}

\columnbreak

\begin{tippbox}
\footnotesize\textbf{Eselsbrücke:}\\[2pt]
\textbf{hay} reimt auf \textit{dabei} → „Was ist alles dabei?"\\
\textbf{ESTAR} = \textbf{E}xakter \textbf{STA}ndo\textbf{R}t \de{(Wo genau?)}
\end{tippbox}

\end{multicols}

\abmark{4, 5}

% ═══════════════════════════════════════════════════════════════
% MERKKASTEN
% ═══════════════════════════════════════════════════════════════

\vspace{6pt}

\begin{merkbox}
\textbf{Zusammenfassung: hay vs estar}

\vspace{4pt}

\begin{multicols}{2}

{\small
\textbf{HAY} -- Existenz (es gibt)
\begin{itemize}[leftmargin=1.2em, itemsep=1pt]
    \item + un/una: \es{Hay una mesa.}
    \item + Zahl: \es{Hay tres libros.}
    \item + mucho/poco: \es{Hay muchas sillas.}
    \item + ohne Artikel: \es{Hay libros.}
\end{itemize}}

\columnbreak

{\small
\textbf{ESTAR} -- Ort (wo ist?)
\begin{itemize}[leftmargin=1.2em, itemsep=1pt]
    \item + el/la: \es{El libro está aquí.}
    \item + mi/tu: \es{Mi mochila está allí.}
    \item + Präposition: \es{está en/encima de...}
\end{itemize}

\vspace{3pt}

\textit{Frage:} \es{¿Dónde está...?}}

\end{multicols}
\end{merkbox}

\end{document}
